% ============================================================================
% PLANTILLA BASE BEAMER - PRESENTACIONES BOOKMATE
% Curso: CC341 - Ingeniería de Software
% Universidad Nacional de Ingeniería (UNI)
% ============================================================================

\documentclass[aspectratio=169]{beamer}
\usetheme{Madrid}
\usecolortheme{default}

% Configuración de fuente
\usepackage[utf8]{inputenc}
\usepackage[spanish]{babel}
\usepackage{graphicx}

% Información de portada
\title{BookMate}
\subtitle{Sistema de Recomendación Inteligente de Libros}
\author{Grupo 6.2}
\institute{Universidad Nacional de Ingeniería\\CC341 - Ingeniería de Software}
\date{\today}

\begin{document}

% Portada
\frame{\titlepage}

% Índice
\begin{frame}{Contenido}
    \tableofcontents
\end{frame}

% Sección 1
\section{Introducción}
\begin{frame}{Introducción}
    \begin{itemize}
        \item Problema: Descubrimiento de libros
        \item Solución: BookMate
        \item Tecnologías: Flask, PostgreSQL, IA
    \end{itemize}
\end{frame}

% Sección 2
\section{Arquitectura}
\begin{frame}{Arquitectura del Sistema}
    \begin{figure}
        \centering
        \includegraphics[width=0.7\textwidth]{ejemplo_arquitectura.png}
        \caption{Diagrama C4 - Contenedores}
    \end{figure}
\end{frame}

% Sección 3
\section{Demo}
\begin{frame}{Demo en Vivo}
    \begin{center}
        \Huge Demo del Sistema
    \end{center}
\end{frame}

% Cierre
\begin{frame}{Conclusiones}
    \begin{itemize}
        \item Objetivos alcanzados
        \item Aprendizajes clave
        \item Trabajo futuro
    \end{itemize}
\end{frame}

\end{document}

